\section{Implementation order and falsifiable next milestones}
\subsection{First: one end-to-end unbounded safety theorem}
Implement the Petri model, finite safety checker, and its generic induction theorem. Use the supplied mutex certificate as a readable first fixture, but do not stop at proving an abstract net safe. Define a small resource program and prove the source simulation. The acceptance criterion is a theorem about that original program, with no external solver required during recompilation and with its actual axiom dependencies printed and reviewed.

The intentionally false sibling should permit two owners to acquire the same lock. Its extracted net should produce a concrete violating word, and the source replay should exhibit the violation. A false sibling is important here: a bridge that accidentally drops the acquire operation or misorders the resource coordinates can make every program appear safe while the abstract checker remains correct.

After this works, add nonconservative nets that defeat a finite forward enumeration, larger randomly generated nets, and independent SMPT comparisons. Keep bounded-oracle exhaustion separate from agreement. Search optimizations come after source semantics and certificate replay, not before.

\subsection{Second: a weighted-step enclosure theorem}
Prove the nonnegative directional remainder lemma directly for sparse monomials. Prove the finite maximum principle and ledger induction. Start with \eqref{eq:cubic}; the resulting Lean theorem should say that the \emph{least} fixed point lies in the checked interval. A proof of $P(r)=r$ alone is not the acceptance criterion.

The first negative fixtures should be the wrong fixed-point branch, an upward-rounded lower step that crosses $q$, an invalid weight, and equality $Jw=w$ falsely used as a strict lower-step witness. Replaying the supplied JSON certificates through a Lean decoder then becomes a useful integration test, rather than the only basis for believing the mathematics.

Only after this abstraction-level theorem works should a branching DSL bridge prove that the interval contains a source termination probability. This sequencing prevents a large probability-library integration from obscuring a small algebraic checker defect.

\subsection{Third: exact critical and algebraic conclusions}
The critical certificate needs finite graph connectivity and the product-deficit inequality in addition to the weighted comparison machinery. Its first pair of acceptance tests is $P(x)=(1+x^2)/2$ versus $P(x)=x$. The former must yield probability one; the latter must not. This pair tests the exact condition that distinguishes branching from endless one-child continuation.

The algebraic decoder should then identify \eqref{eq:cubic-root} using factorization and an isolating interval. Its first extension can reuse the baseline's existing scalar sign-certificate plans for factors with repeated roots or difficult derivative intervals. A new general algebraic-number engine is not required to deliver the first exact irrational result.

\subsection{Fourth: synthesis and larger semantic fragments}
Once source bridges exist, wire them to retained Leant/Djex candidate graphs for a small provider inventory. Measure source-contract acceptance, counterexample replay, and undecided cases separately. Preserve the original synthesis query, provider identities, and limits; replacing a hard original query by a simpler hand-written program is not acceptance of that query.

For probability systems, implement strongly connected component processing with exact-value substitution first, and interval-parameter transport second. For multiset systems, add dominance indexing and certificate minimization. For Newton search, use sparse exact solves and numerical proposals followed by rational residual checks. A numerical producer can propose $d,w$ directly for \eqref{eq:weighted-step}; it does not have to reconstruct an entire inverse.

Each step has an explicit mathematical contract and an available finite checker target. This makes a failure actionable: unsupported source semantics, no certified lower progress, no prefixed upper candidate, insufficient precision, or resource exhaustion are distinct engineering outcomes.

\section{Assessment}
The strongest extension is not another collection of local tactics. It is the ability to turn two genuinely infinite questions into small finite obligations of the right logical polarity.

For multiset-state safety, a finite backward basis proves a forward invariant without enumerating all reachable markings. For recursive probabilities, lower histories and prefixed upper vectors distinguish a least solution from arbitrary roots. A weighted maximum principle removes full matrix inverses from many Newton certificates. At criticality, an exact nondegeneracy theorem proves probability one where strict contraction is unavailable. A factor identity and rational interval then make an irrational probability an exact mathematical result rather than a decimal guess.

The delivered package demonstrates those mathematical building blocks with exact arithmetic, independent finite-oracle comparisons where available, explicit negative controls, search-disabled replay, and deterministic reproduction. It does not yet close the source-to-Lean proof chain. The module plan and acceptance criteria identify that remaining work precisely, without claiming that another Python experiment could substitute for it.
